\section{完整案例：等间隔趋势汇总}
同表采用固定纵向范围，不逐行放大微小波动；缺失月份保留位置。
\begin{example}[text and listing,sidebyside=false]{示例 73：半年观测与末期数值}
\gradbarsstyle{monthly}{width=65mm,height=7mm,
  min=0,max=100,spark range=fixed,precision=0,
  row padding=3pt,spark points=extrema}
\begin{center}
\begin{tabular}{@{}cc@{}}
\toprule
\multicolumn{1}{c}{Series} & \multicolumn{1}{c}{Jan--Jun / June value}\\\midrule
North & \gradspark[style=monthly]{35,48,42,60,72,80}\\
South & \gradspark[style=monthly]{55,53,NA,58,62,65}\\
West & \gradspark[style=monthly]{45,45,45,45,45,45}\\
East & \gradspark[style=monthly]{30,40,55,60,70,NA}\\
\bottomrule
\end{tabular}
\end{center}
\end{example}
North 与 South 的极值点只表示各自的最高与最低观测，不表示好坏。
East 的末期标签为缺失，不能用前一期的 70 代替。
不等间隔日期应先整理到公共观测网格；不要直接当作等间距序列。
